\subsection{How to Tell Whether a Milestone Has Been Met}
\label{sec:whether-milestone}
A milestone claimed by the organization that built the system is a different kind of evidence than one checked by somebody else, and the gap is empirical rather than a matter of trust. A system's own claim inherits the limits on self-report: fluent confidence does not establish correspondence with its internal states, and the concept-injection entry (§\ref{sec:A.selfreport}) puts a rate on the gap. A passing score on a fixed test suite fails differently, inheriting specification gaming and goal misgeneralization. What either needs instead is a passing score under conditions built to break it.

Who runs that harder test matters as much as the test itself. The limit on self-regulation applies here without much modification: a review conducted or funded by the organization being reviewed answers a structurally different question than one conducted by a party with no stake in the answer, whatever the intention behind the first. METR, a nonprofit that evaluates frontier AI systems' capabilities, runs assessments both in partnership with developers like Anthropic and OpenAI and, separately, on its own after a system has already been released \autocite{metr2026about}. The United Kingdom's AI Security Institute does government-run evaluation of frontier models before deployment and is an independent evaluator, not a rubber stamp — but by its own account it is not a regulator: it cannot compel a company to submit a model for testing, block a release, or intervene once a system is already causing harm \autocite{aisi2025governance}. Independent evaluation and enforcement are not the same capability, and a milestone checked by the first without access to the second is still short of the binding external answer self-regulation structurally cannot provide on its own.

Independence is not competence: past some capability level the rater's verdict stops being reliable evidence of the property scored, and an evaluation in that region has to name the limit rather than report a number that looks more settled than the method is.

What an evaluation like this is for is not confirming a milestone but being willing to report that it wasn't met — and, where the milestone turns out to have measured the wrong thing, revising the milestone. A negative result may mean that the work fell short or that the milestone measured the wrong thing. Evaluators need to distinguish those explanations before deciding what to revise. 